\scanpage{7}
Exemples: Si $S$ est la catégorie des schémas affines, avec une topologie plus fine que Zar, $M$ l'ensemble des immersions ouvertes, \deleted{alors la const}\deleted{la construction donn\unsure{é|e}e} (ou, en l'absence de d), les immersions de la forme $\mathrm{Spec}(A_f)\handwritten{\to}\mathrm{Spec}(A)$), alors la construction donne une catégorie $\overset{\handwritten{\sim}}{S}_M$ équivalente à celle des schémas. Si on prend une topologie plus fine que la topologie étale, et pour $M$ \inserted{\deleted{\unsure{h}}l'ensemble} des morphismes étales (qui satisf\deleted{\unsure{on}}ait la condition d)), alors on trouve les ``espaces algébriques'' de Artin, en ``recollant des spectres d'anneaux \deleted{pour \unsure{l}} par des morphismes étales''. Le dernie\deleted{\unsure{f}}r corollair\unsure{e|è} nous dit qu'on aurait trouvé le même résultat en partant de la catégorie $S'$ des schémas, au lieu de celle $S$ des schémas affines. \handwritten{Dans ce cas, les morphismes de $\widetilde M$ sont dits \uline{étales} [la top. sur $S$ étant la top. étale). \deleted{\illegible{1}}}

\uline{Proposition} \quad Supposons que $M$ satisfasse la condition draconienne suivante: soient $u:X\to Y$ et $v:Y\to Z$ des flèches composables dans $S$, \deleted{\unsure{alor|alors}} alors $v\handwritten{\in} M$ et $vu\handwritten{\in} M$ implique $u\handwritten{\in} M$. Pour un objet $X$ de $S$, soit $S(X,M)$ la sous-catégorie de $S_{\handwritten{/}X}$ \deleted{formé} formée des $X'\to X$ qui sont dans $M$ (donc la condition précédente dit que tous les morphismes de $S(X,M)$ sont dans $M$); munissons-là de la topologie induite par $S$, et considérons le topos $S(X,M)^{\handwritten{\sim}}$. Alors \deleted{le foncteur $(\overset{\handwritten{\sim}}{S}_M)_{\handwritten{/}X}\quad S(X,M)^{\handwritten{\sim}}$ composé de}
\diagram{\begin{tikzcd}[ampersand replacement=\&,column sep=large]
\deleted{(\overset{\handwritten{\sim}}{S}_{M\unsure{x}})_{\handwritten{/}X}}
\arrow[r,"{\deleted{\mathrm{incl.\ can.}}}"] \&
\deleted{\overset{\handwritten{\sim}}{S}_{\handwritten{/}X}}\arrow[r,"{\deleted{\handwritten{\approx}}}"] \&
\deleted{(S_{\handwritten{/}X})^{\handwritten{\sim}}}\arrow[r,"{\deleted{\mathrm{restr}}}"] \&
\deleted{S(M,X)^{\handwritten{\sim}}}
\end{tikzcd}}
\deleted{est une équivalence de catégories.} la sous-catégorie pleine de $\overset{\handwritten{\sim}}{S}_{\handwritten{/}X}$ formée des $F\handwritten{\to}X$ qui sont dans $\overset{\handwritten{\sim}}{M}$ \inserted{\handwritten{[d'où résulte $F\in\mathrm{Ob}\,\widetilde S_M$ !]}}, est \deleted{équivalent} un \uline{topos}, équivalent par le foncteur restriction à $S(\deleted{M}X,\deleted{X}M)$ au topos $S(X,M)^{\handwritten{\sim}}$.

% archive sheet number: 6
